\section{Galois representations of lifts}\label{sec:galois-lifts}

Part~\ref{part:lower} reduced Theorem~\ref{thm:C} to the existence of a regular family of twists (\cref{thm:regular-family}): for each prime place $\iota$ over $\ell$, an $m$-fold twist
$\xi_\iota$ whose powers have few congruences, in the sense that the denominators of the projections to the lifts are small, and whose inverse ratios of Hecke values are
barely divisible by $\ell$. This part constructs such twists. Congruences between a lift and forms orthogonal to the lifts produce, by Ribet's method, Selmer classes
in the Galois representations attached to the lifts (\cref{sec:ribet,sec:denominators}); the local conditions above $\ell$ cost a bounded loss
(\cref{sec:local-loss,sec:crystalline}); and a twist ramified at new auxiliary primes makes all these Selmer groups short and keeps the ratios of Hecke values
barely divisible (\cref{sec:selmer,sec:twisting}). \Cref{sec:fit} assembles the pieces. That denominators of projections to theta lifts give Selmer classes, and that twisting lowers Selmer ranks, is known in other settings \cite{Prasanna2006,Klosin2009,MazurRubin2010}; here the bounds are needed for all primes $\ell$ at once, summed with weights $\log\ell$.

This section describes the Galois representations that occur: those attached to the eigenforms in a space of lifts are, on traces, a twisted \emph{main summand},
independent of the twist, plus a string of powers of the inverse cyclotomic character.

\subsection{Conjugate self-dual representations}\label{ss:conj-self-dual}
Fix a prime $\ell$ and an isomorphism $\iota\colon\C\simeq\bar\Q_\ell$. Let $G=\Gal(\Omega/E)$ for the field $E$ of the data, $c$ the complex conjugation, $g^c=cgc$, and $\mathrm{cyc}$ the $\ell$-adic cyclotomic
character. The \emph{bad places} are the places of $E$ dividing $\ell p_0p_1N$. The \emph{avatar} $\vartheta^{-1}_\ell$ of a tame character $\vartheta$ is the $\ell$-adic Galois character of the base change of $\vartheta^{-1}$ to $E$,
through $\iota$. A continuous representation $R\colon G\to\GL_n(\bar\Q_\ell)$ is \emph{conjugate self-dual of twist $1-t$} if its trace form on the span of its image is nondegenerate and
$\operatorname{tr}R(g^c)=\mathrm{cyc}(g)^{1-t}\operatorname{tr}R(g^{-1})$ for all $g$; it is \emph{of level $S$} with jumps $J$ if moreover, at every place $v$ outside $S$ or above $\ell p_0p_1$, it is trivial on inertia near $1$
(potentially crystalline above $\ell$), has the Hodge--Tate jumps $J(\sigma)$ at the embeddings $\sigma$ over $v$, and is unramified (crystalline above $\ell$) if $v\nmid p_0p_1$.
\lean{New.Data.algebraMap_mem_cycloField, New.Selmer.AbsGal, New.Selmer.placesAboveM, New.Selmer.badPlaces, New.Selmer.placesDividing_union_subset_badPlaces, New.Data.IsChosen.mem_placesAboveM, New.Selmer.star_mem_placesAboveM, New.Selmer.norm_natCast_le_norm_sub_of_pow_eq_one, New.Selmer.continuous_cyclotomicInv, New.Selmer.cyclotomicInv_frob, New.Data.TameChar.invAvatar, New.Selmer.avatar_conjGal_mul_frob, New.Selmer.avatar_conjGal_mul_avatar, New.Selmer.isLabelledAt_avatar, New.Data.TameChar.tamePart_eq_one_of_isChosen, New.Data.isChosen_levelPlaces_chosenPlace, New.Selmer.isChosen_chosenPlace, New.Selmer.IsConjSelfDual, New.Selmer.twistError, New.Selmer.selmerCocycles, New.Selmer.IsTwistedLattice, New.Selmer.IsTwistedLattice.cyclotomicChar_eq_one, New.Selmer.HasTwistedLattice, New.Selmer.crystallineSelmerLength, New.Selmer.HasDivisibleSelmerLengthLe, New.Selmer.length_selmerCocycles_le_crystallineSelmerLength, New.Selmer.HasDivisibleSelmerLengthLe.hasTwistedLattice, New.Selmer.HasDivisibleSelmerLengthLe.mono, New.Selmer.IsConjSelfDualOfLevel, New.Selmer.IsMainSummandLike, New.Selmer.IsMainSummandLike.eventually_span_eq_top_and_exists_eq_smul_one, New.Selmer.IsMainSummandLike.eventually_eq_bot_or_eq_top, New.Selmer.IsMainSummandLike.eventually_lattice_unique, New.Selmer.normalizingAvatar, New.Selmer.normalizingAvatar_eq_of_isFrob_inv, New.Selmer.normalisedRep, New.Selmer.scaledTypePoly, New.Selmer.IsLiftSetting, New.Selmer.typeProjection, New.Selmer.HasLiftTrace, New.Selmer.IsMainSummand, New.Selmer.IsRepOfType, New.Selmer.IsLiftSetting.card_le, New.Selmer.IsLiftSetting.isChosen_chosenPlace, New.Selmer.IsLiftSetting.isSymmetric_typeProjection, New.Selmer.IsLiftSetting.typeProjection_eq_self_of_mem_range, New.Selmer.IsLiftSetting.range_typeProjection_mem_invtSubmodule, New.Selmer.frequently_isSplit_notMem_places}
We record the elementary facts about these notions that the later sections use. For a character $a$ of $[\U(1)]$ of weight $e$ and tame part $\psi$ trivial off the chosen places, and integers $k$, $j$, the
\emph{avatar} of the algebraic Hecke character $y\mapsto a(y/\bar y)\upsilon(y)^{-k}\abs y^{j-k/2}$ is its $\ell$-adic Galois character through $\iota$ (\cref{thm:B-galois}(a) in rank one); $\upsilon$ is the splitting character
(\cref{thm:B-splitting-character}). A matrix representation $u$ over the ring of integers $\OO$ of a finite extension of $\Q_\ell$ is a \emph{lattice} in $R$ if $R$ is conjugate over $\bar\Q_\ell$ to $u$ read in $\bar\Q_\ell$. A
\emph{coefficient ring} is the ring of integers of a finite Galois extension of $\Q_\ell$, and \emph{all large} coefficient rings are those containing a given one.

\begin{lemma}[Bad places, avatars and lattices]\label{lem:avatars}
\begin{enumerate}[label=\textup{(\alph*)}]
\item The places above $\ell$ and above $p_0p_1$ are bad. A chosen place lies above $p_0$ or $p_1$, and so does its conjugate; the tame part of a tame character is trivial at the chosen places. Every set of places
above $p_0p_1$ has at most $4f$ elements, and some chosen place not above $\ell$ has residue field of order $p_0$ or $p_1$.
\item At a Frobenius element $g$ at a place $w\nmid\ell$ the inverse cyclotomic character takes the value $q_w$. Two distinct $N$-th roots of unity in $\bar\Q_\ell$ have distance at least $\abs N_\ell$.
\item The avatar $\mu$ of $a(y/\bar y)\upsilon(y)^{-k}\abs y^{j-k/2}$ satisfies $\mu(g^c)\mu(g)=\mathrm{cyc}(g)^{2j-k}$ for all $g$. At a split place $w$ where $\psi$ is trivial at $\bar w$, inertia acts on $\mu$ through the element of the
residue field that names $\psi_w$: an element of tame exponent $i$ is killed by $X-\iota(T^i(\text{name}))$, $T$ the faithful character of the algebraic closure of the residue field.
\item A lattice $u$ in $R$ has the traces and determinants of $R$, is continuous if $R$ is, and is trivial where $R$ is.
\end{enumerate}
\end{lemma}

\begin{proof}
(a) A prime ideal containing one factor of $\ell p_0p_1N$ contains the product. The chosen places divide $\pi_M$, an integer of $K$ with $\pi_M\bar\pi_M=p_0p_1$, so they lie above $p_0$ or $p_1$; the conjugates of places
above $p_i$ lie above $p_i$. The places of the modulus of a tame character are not chosen, and its tame part is trivial off its modulus. Over each $p_i$ there are $f$ chosen places, and a place above $p_i$ is chosen or the
conjugate of a chosen one, so there are at most $2\cdot2f$ places above $p_0p_1$. A place above both primes would have residue fields of two characteristics; so for $p_i\ne\ell$ a chosen place above $p_i$ is not above $\ell$,
and its residue field has $p_i$ elements: $q_1d\mid p_i-1$ (\cref{prop:sizes}) and $E\subseteq\Q(\zeta_{q_1d})$ (\cref{prop:field-E}(a)), so $p_i$ splits completely in $E$.

(b) The cyclotomic character at a geometric Frobenius at $w\nmid\ell$ is $q_w^{-1}$. If $\zeta\ne\zeta'$ are $N$-th roots of unity, $\prod_{\eta\ne1}(1-\eta)=N$ over the $N$-th roots of unity $\eta$, and every factor has absolute value at
most $1$, so $\abs{1-\zeta'/\zeta}\ge\abs N$.

(c) Both sides are continuous characters, so by Chebotarev density (\cref{thm:B-classfield}(d)) it suffices to check them at Frobenius elements at places $w\nmid\ell$ where $\psi$ is trivial at $w$ and $\bar w$. There
$g^c$ is a Frobenius element at $\bar w$, the avatar at a Frobenius element is $\iota$ of the value of the Hecke character at a uniformiser, and in the product of the values at $\bar w$ and at $w$ the factors of the base
change of $a$ and of $\upsilon$ cancel ($a$ is a character of $\U(1)$ and $\upsilon(\bar y)=\upsilon(y)^{-1}$), leaving $\abs{\varpi_w\bar\varpi_w}^{j-k/2}$ at the two places, which is $q_w^{-(2j-k)}=\mathrm{cyc}(g)^{2j-k}$ by (b). For the
inertia statement, an element $s$ of tame exponent $y$ is the Artin image of a local unit $x$ (local class field theory with the roots of unity of order $5040$ in $E$, \cref{thm:B-classfield}), so it acts by the value of the
character at $x$, which is $\psi_w$ at the residue of $x$; this value is $T^y$ of the name of $\psi_w$.

(d) Some injective $\OO$-algebra map from the matrices over $\OO$ to the matrices over $\bar\Q_\ell$, an inclusion of the entries followed by a conjugation, carries $u(g)$ to $R(g)$ and preserves traces and determinants; it is
continuous and injective.
\end{proof}


\begin{definition}[Main summand]\label{def:main-summand}
Let $1\le r\le\mathrm{r}$ (the rank of the data) and $t\le\mathrm t=5\mathrm r$. Jumps $J(\sigma)$, multisets of integers indexed by the complex embeddings, are \emph{main summand weights} for $(r,t)$ if, for one $j\le m$, every jump $u$
at $\sigma$ is \emph{low}, $\Delta/4\le\abs u\le2m\Delta$, or \emph{high}, $\abs{\epsilon_\sigma u+w_j(\sigma)}\le\mathrm t$ with $\epsilon_\sigma=\pm1$ the sign of $\sigma$ relative to the CM type and $w_j$ the weight of the $j$-th lift; at
most $rm/k$ jumps at $\sigma$ have $\abs u\le2(\Delta/2)^2$; and at most $r$ jumps at $\sigma$ and $\bar\sigma$ together are negative. A representation $R$ of rank $r$ is \emph{like a main summand} if it is conjugate self-dual of
level $\emptyset$ with main summand weights and, in addition: at the places above $p_0p_1$ prime to $\ell$ a power prime to the place kills the image of inertia, and for one inertia element $\gamma$ the determinants
$\det(R\gamma-1)$ and $\det(R(\gamma^c)^{-1}-1)$ have valuation at most $r\,v_\ell(n_{\rm ch})$, $n_{\rm ch}$ the order of $\prod_w\GL_{\mathrm r\mathrm t}(\F_w)$ over the chosen places; over all large coefficient rings the image of a
lattice spans all matrices and contains a scalar at an element of cyclotomic character $p_i^r$; and some inertia element above $p_0p_1$ acts by a scalar $y$ with $y^{\ell^aj}-1$ a unit for all $a$ and $1\le j\le4r+2$.
\end{definition}

\begin{proposition}[Main summands]\label{prop:main-summands}
A representation like a main summand is residually absolutely irreducible over all large coefficient rings, and its stable lattice is unique up to scalars.
\end{proposition}

\begin{proof}
\emph{Irreducibility.} Let $\OO$ be large, $u$ a lattice in $R$ over $\OO$ and $f\colon\OO\to k$ a ring map to a field. The image of $u$ spans all matrices over $\OO$ (\cref{def:main-summand}), so the images $f(u(g))$ span all matrices
over $k$, and a subspace of $k^n$ stable under them is stable under all matrices, hence $0$ or $k^n$. \emph{Uniqueness of the lattice.} Let $u$, $u'$ be the matrices of $R$ over $\OO$ in the bases $Q$, $Q'$, and $B=Q^{-1}Q'$, so that
$Bu'(g)B^{-1}=u(g)$. Multiply $B$ by a scalar so that $L=B\OO^n\subseteq\OO^n$ and $L\not\subseteq\varpi\OO^n$, $\varpi$ a uniformiser. Then $L$ is stable under the $u(g)$, so its image in $(\OO/\varpi)^n$ is a nonzero stable subspace,
which by irreducibility is everything; by Nakayama's lemma $L=\OO^n$, so the scaled $B$ lies in $\GL_n(\OO)$: $Q'=\alpha Qk$ with a scalar $\alpha$ and $k\in\GL_n(\OO)$.
\end{proof}

\subsection{Representations of a type}\label{ss:rep-type}
A descent pair $O$ with labels $c=(c_i)$ at its chosen places is a \emph{lift setting} for a twist $\vartheta$ if its places are those of the data and of $\vartheta$, its tame part is $\omega^j\vartheta$ with $1\le j\le m$, the
rank $r_O$ and the difference of ranks $t_O$ are odd with $r_O\le\mathrm r$, $t_O\le\mathrm t$, each $c_i$ is a label of degree $r_O$ whose \emph{scaled} polynomial (roots multiplied by the name of $\psi_w^{-1}$) does not have the root
$1$, and over each prime $p_k$ one label is primitive with $r_O=\mathrm r$, or is $X-1$ with $r_O=1$. A representation is \emph{of type $c$} if it is conjugate self-dual of a given twist, jumps and level, and inertia at
the chosen places acts through the roots of the scaled labels. The \emph{normalised representation} of $R$ is its twist by the avatar of the inverse of the character of $W$, of weight $-r_O$.
\lean{New.Selmer.not_isRoot_one_scaleRoots_inv, New.Selmer.isLiftSetting_descentPair, New.Selmer.isLiftSetting_mainPair, New.Selmer.trivialLift_mem_range_typeProjection, New.Selmer.levelPrime_pow_sub_one_dvd_chosenGroupOrder, New.Selmer.ellVal_levelPrime_pow_sub_one_le, New.Selmer.card_le_of_subset_placesAboveM, New.Selmer.exists_chosen_notMem_placesDividing, New.Selmer.IsLatticeIn.exists_algHom, New.Selmer.IsLatticeIn.continuous, New.Selmer.IsLatticeIn.eq_one, New.Selmer.IsLatticeIn.trace_eq, New.Selmer.coeff_zero_mul_coeff_star, New.Selmer.isConjSelfDual_of_attached, New.Selmer.isTraceNondeg_transpose_inv, New.Selmer.count_jumps_comp_star, New.Selmer.isFiniteTypeAt_and_hasJumpsAt_of_conjPlace, New.Selmer.prod_aroots_dvd_prod_aroots_scaleRoots_comp, New.Selmer.IsRepOfType.twist_avatar, New.Selmer.isRepOfType_of_hasProductLabel}
\lean{New.Selmer.hasParahoricType_chosenGLRep, New.Selmer.coe_asAlgebraHom_chosenGLRep_apply, New.Selmer.componentOn_places_eq_self, New.Selmer.toSmoothForm_mem_levelSpace_componentOn, New.Selmer.toSmoothForm_typeProjection, New.Selmer.chosenGroupOrder_nsmul_typeProjection_mem_ellIntegralForms, New.Selmer.notMem_and_star_notMem_range_union_places, New.Selmer.normalisedJumps, New.Selmer.normalisedJumps_eq_map_hcParameter, New.Selmer.normalisedJumps_eq_range_add_map, New.Selmer.apply_ofSplitPlace_eq_formTranslation, New.Selmer.hasProductLabel_of_tupleIdem_apply_ne_zero, New.Selmer.hasTwistParahoricType_levelSpace, New.Selmer.toSmoothForm_mem_levelSpace_repOff, New.Selmer.exists_mem_levelSpace_chosenGLRep_map_eq, New.Selmer.translationAtPlaces, New.Selmer.IsFormOfType, New.Selmer.IsFormOfType.isRepOfType_normalisedRep, New.Selmer.exists_baseLevel, New.Selmer.isRepOfType_normalisedRep, New.Selmer.IsFormOfType.isRepOfType_normalisedRep}
We use three facts about lift settings and types. A \emph{good tuple} of labels at the chosen places is a family of labels of degree $\mathrm r$ (monic, squarefree, nonzero at $0$) vanishing at no $q_2$-th root of unity, $q_2$ the
order of the residue characters $\omega$ at the chosen places, with one primitive label over each of $p_0$, $p_1$. The \emph{type projection} $e_c$ is the action on forms of the idempotent of the tuple $c$ in the group algebra of
$\prod_{w\in M}\GL_{r_O}(\F_w)$. The normalised jumps of $O$ are the jumps of the weight of the forms on $V$ minus the infinity type of the normalising character.

\begin{lemma}[Lift settings]\label{lem:lift-settings}
\begin{enumerate}[label=\textup{(\alph*)}]
\item A good tuple $c$ forms a lift setting with the $j$-th descent pair, for every twist $\vartheta$ and $1\le j\le m$. For a main pair $p$ for $j$, the descent pair of the lift of $\nu\vartheta$ from $\U(1)$, with every label $X-1$, is a
lift setting for $\vartheta$, and the lift $\varphi_{X_{p,\vartheta}}$ lies in the range of its type projection.
\item For $1\le d\le\mathrm r\mathrm t$, $p_i^d-1$ divides $n_{\rm ch}$, so $v_\ell(p_i^d-1)\le v_\ell(n_{\rm ch})$.
\item For a lift setting, $e_c$ is self-adjoint and fixes its range, which is stable under the Hecke polynomials at the split places off the level; infinitely many split places are, with their conjugates, off the level.
\item Let $f$ be an eigenform of weight $\kappa$ on the definite unitary group of rank $n'$, $R$ attached to it, and let $f$ be spherical at the places off $S$, above $\ell$ or above $p_0p_1$, except the chosen places and their
conjugates, with a type at each chosen place induced from the label $c_i$ of degree $r$ and trivial characters on blocks. Then $R$ is of a type: conjugate self-dual of twist $1-n'$ and level $S$ with the jumps of $\kappa$, and
inertia at the chosen places acts through the roots of $c_i(X-1)^{n'-r}$.
\item Twisting a representation of a type by the avatar of $a(y/\bar y)\upsilon(y)^{-k'}\abs y^{j'-k'/2}$, with the tame part of $a$ trivial off the chosen places, gives a representation of a type of the same level, with twist
shifted by $2j'-k'$, jumps lowered by the infinity type, and the roots of the labels multiplied by the element naming the tame part of $a$ at the chosen place.
\item If $n'$ is even and $l=2\lambda$ fits $O$, the normalised jumps at $\sigma$ are $0,1,\dots,t_O-1$ and $\frac{t_O-1}2\mp(\lambda_c+\kappa')$, $c$ running over the basis of $W$, with the upper sign on the CM type; $-\kappa'$ is the weight of
the character of $O$.
\end{enumerate}
\end{lemma}

\begin{proof}
(a) \emph{Descent pairs.} The places of the $j$-th descent pair are those of the data and of $\vartheta$, its tame part is $\omega^j\vartheta$, and its ranks $\mathrm r$ and $\mathrm t$ are odd (\cref{prop:sizes}). The scaled label at $w$ has
the root $1$ only if a root of $c_w$ names $\omega^j$ at $w$, a root of unity of order dividing $q_2$, which a good label does not have; if $y$ is a unit that is not a root of $c$, then $1$ is not a root of $c$ with its roots
multiplied by $y^{-1}$. \emph{The lift from $\U(1)$.} There $W$ has rank $1$ and the difference of ranks is $2z-1$, both odd, and $1\le\mathrm r$, $2z-1\le\mathrm t$. The scaled label of $X-1$ has the root $1$ only if the tame part of $\nu$
is trivial at the place, and a main pair is ramified (\cref{def:critical-pair}, \cref{ss:twist-families}). The label $X-1$ belongs to the trivial character of $\GL_1$, so the type projection is the mean over
$\prod_{w\in M}\GL_1(\F_w)$, which fixes $\varphi_{X_{p,\vartheta}}$ (\cref{prop:norm-L-value}(b)).

(b) $p_i^d-1$ is the first factor of $\#\GL_d(\F_{p_i})=\prod_{e<d}(p_i^d-p_i^e)$, which divides $\#\GL_{\mathrm r\mathrm t}(\F_w)$ at a chosen place $w$ above $p_i$, and $n_{\rm ch}$ is a multiple of these orders.

(c) The idempotent is a central idempotent of the group algebra, hence self-adjoint (\cref{lem:hecke-adjoint}(b)) and the identity on its range. A group element acts by a root of unity times translation by a point that is $1$
at $w$ and $\bar w$, which commutes with the translations supported at $w$ and $\bar w$, hence with the Hecke polynomial. Places of prime norm split, there are infinitely many, and finitely many lie in the level or its conjugate.

(d) Off the chosen places this is \cref{thm:B-galois}(a) together with conjugate self-duality, proved in \cref{thm:attached}(a) below without using (d). At a place over a chosen place $w$ the type of $f$ is a product label at $w$,
so \cref{thm:B-galois}(c) gives finite type, the jumps and a polynomial built from the tame characters of the label, $c_i$ for the block of size $r$ and $X-1$ for the trivial blocks of total size $n'-r$, that kills inertia.
At the conjugate places these pass by conjugate duality: $R\circ(g\mapsto g^c)$ is conjugate to the twisted dual of $R$, both having nondegenerate trace forms and the same traces, conjugating an embedding replaces a jump $u$
by $n'-1-u$, and a polynomial killing $s$ at $\bar A$ kills $(s^c)^{-1}$ at $A$.

(e) Each local clause passes to the twist from the clause for $R$ and that for the character. The avatar is of finite type with the single jump minus the infinity type; off $p_0p_1$ the character is unramified,
neither the place nor its conjugate being chosen; and an inertia element $s$ acts on the avatar by $\varphi(u)$, $u$ the name of the tame part, and on $R$ through the values of $\varphi$ at the roots of $c_i$, so on the twist
through these values times $\varphi(u)$, as $\varphi(u)X-\varphi(ux)=\varphi(u)(X-\varphi(x))$. The twist of the cyclotomic exponent is \cref{lem:avatars}(c).

(f) The normalising character has exponent $r_O-r_-$ on the CM type and $r_-$ off it, $r_-$ the number of minus signs of $W$, and entry $p$ of twice the Harish-Chandra parameter is $2\mu_p+n'-1-2p$ for the highest weight
$\mu$. Hence the normalised jumps are $(t_O-1\mp(a_p-2r_-+r_O))/2$, $p<n'$, $a$ the doubled parameter. By the block rule (\cref{prop:descent-pairs}(b)) each entry of $a$ is a shifted entry of $l$ or a member of the doubled block of
length $t_O$, which gives the stated values; both sides are sets of $n'$ distinct integers, so the inclusion is an equality.
\end{proof}


\begin{theorem}[Attached representations]\label{thm:attached}
Let $O$, $c$ be a lift setting for $\vartheta$ with $Q(\vartheta)$ auxiliary (its places of prime norm and, with their conjugates, not bad).
\begin{enumerate}[label=\textup{(\alph*)}]
\item A Galois representation attached to an eigenform on a definite unitary group of rank $n$ is conjugate self-dual of twist $1-n$.
\item For an eigenform $f$ of type $c$ on the larger space of $O$, with the level of $O$ off a set $S$ of places outside the bad places, the normalised representation of an attached $R$ is of type $c$: conjugate
self-dual with the jumps $0,1,\dots,t_O-1$ and $\frac{t_O-1}2\mp(\lambda_c+\kappa')$ (the upper sign on the CM type, $\lambda$ the parameter on $W$ and $-\kappa'$ the weight of the character of $O$), of level $Q$, $S$
and their conjugates, with inertia at the chosen places acting through the roots of the scaled labels times $(X-1)^{t_O}$.
\item There are base data whose forms contain the images of the type projection $e_c$ of the forms of $O$; $n_{\rm ch}e_c$ keeps integral forms; and the normalised representations of their eigensystems
of auxiliary level are of type $c$, with $(u(g)-\vartheta_\ell^{-1}(g))(u(g)-1)=0$ for $g$ in inertia above $Q$ and its conjugates.
\end{enumerate}
\end{theorem}
\lean{New.Selmer.norm_pow_sub_pow_eq_one, New.Selmer.ellVal_det_sub_one_le_card_mul, New.Selmer.not_q₂_dvd_mul_pow_mul, New.Selmer.IsInertiaSpectrum, New.Selmer.IsInertiaSpectrum.of_conjPlace, New.Selmer.exists_isInertiaSpectrum_chosenPlace, New.Selmer.exists_inertiaGroup_isInertiaSpectrum, New.Selmer.IsRepOfType.exists_map_pow_eq_one_and_ellVal_det_le, New.Selmer.IsCommonRootPower, New.Selmer.IsLiftSetting.exists_isCommonRootPower_of_isPrimitivePolynomial, New.Selmer.IsLiftSetting.exists_isCommonRootPower_of_eq_X_sub_one, New.Selmer.IsLiftSetting.exists_isCommonRootPower, New.Selmer.IsRepOfType.exists_inertia_isRoot_charpoly, New.Selmer.eventually_span_eq_top_and_exists_scalar_of_tame_pair, New.Selmer.IsRepOfType.isMainSummandLike, New.Selmer.invAvatar_eq_of_isTameExponent, New.Selmer.isSplit_twistPlace, New.Selmer.HasTwistParahoricType, New.Selmer.natCast_notMem_of_mem_conjClosure, New.Selmer.isFiniteTypeAt_and_polyKillsAt_of_mem_conjClosure, New.Selmer.normalisedRep_quadratic_on_inertia}

\begin{proof}
(a) \emph{A coefficient identity.} Let $f\ne0$ in rank $n=k+1$ be spherical at a split $w$ and at $\bar w$, with eigenvalues $s_i$ and $u_i$ of the $i$-th Hecke operators there, so that its Hecke polynomials are
$Q_w=\sum_i(-1)^iq^{i(i-1)/2}s_iX^{n-i}$ and $Q_{\bar w}$ likewise (\cref{ss:B-galois}). For $i+i'=n$ the $i$-th operator at $\bar w$ followed by the scalar matrix of a uniformiser at $w$ is the $i'$-th operator at $w$, a point of
$\U(V)$ at $\{w,\bar w\}$ being determined by its component at $w$ (\cref{prop:point-forms}(b)). For $i=0$ the scalar acts by $s_n$; for $i=1$ this gives $s_nu_1=s_{n-1}$, that is $c_0(Q_w)\,c_k(Q_{\bar w})=q^k\,c_1(Q_w)$ for the
coefficients $c_d$ of $X^d$. \emph{Passage to the representation.} For a matrix $A$ of size $n$ the coefficients of degree $0$, $1$ and $k$ of $\det(X-A)$ are $(-1)^n\det A$, $(-1)^{k}\det A\operatorname{tr}A^{-1}$ and $-\operatorname{tr}A$. The map
$g\mapsto g^c$ is a homeomorphism carrying Frobenius elements at $w$ to Frobenius elements at $\bar w$, so $R$ read through it is attached to the system $w\mapsto Q_{\bar w}$. For most Frobenius elements $g$ the place is prime to
$\ell$ and $\mathrm{cyc}(g)^{-1}=q_w$ (\cref{lem:avatars}(b)). Along Frobenius elements tending to an arbitrary $g$, which exist by Chebotarev density (\cref{thm:B-classfield}(d)), the coefficients of $\iota(Q_w)$ and $\iota(Q_{\bar w})$
tend to those of the characteristic polynomials of $R(g)$ and $R(g^c)$, so in the limit
\[
  (-1)^n\det R(g)\cdot\bigl(-\operatorname{tr}R(g^c)\bigr)=\mathrm{cyc}(g)^{-k}(-1)^k\det R(g)\operatorname{tr}R(g^{-1}).
\]
The signs agree and $\det R(g)\ne0$, so $\operatorname{tr}R(g^c)=\mathrm{cyc}(g)^{1-n}\operatorname{tr}R(g^{-1})$ for all $g$.
Continuity and the nondegenerate trace form are part of the definition of an attached representation.

(b) Let $f$ be an eigenform of type $c$ with the level of $O$ off $S$. The type projection fixes $f\ne0$, so the idempotent of each label does not kill it, and at each chosen place $f$ has a product label with two blocks: the block
of size $r_O$ carries the class function of the scaled label, the other block the trivial character (the trivial character is named by $1$, which is not a root of the scaled label). A place that is neither chosen nor in $Q$, $S$
or their conjugates is a place of maximal level, because $Q$ and $S$ avoid the bad places. By \cref{lem:lift-settings}(d) the attached $R$ is of a type, with inertia at the chosen places acting through the roots of the scaled
labels times $(X-1)^{t_O}$ (the tame part of the normalising twist is trivial, so the roots are not scaled again); its normalisation by the avatar of the inverse of the character of $W$ is of a type by
\cref{lem:lift-settings}(e), with the normalised jumps of \cref{lem:lift-settings}(f), the rank $r_O+t_O$ being even.

(c) Let $\pi$ act on forms through the components at the places of $O$ only, and let $\mathcal Y$ be the forms with the level and character of $O$ under $\pi$ that the idempotent of $c$ fixes. These are forms of base data of the weight of
$O$ with no extra level: the components of a point at the places of $O$ and their conjugates change neither the reductions there nor the level character, and the maximal level off these places fixes the forms of $O$.
An operator commuting with $\pi$ on the maximal level commutes with the idempotent, so the Hecke operators off the level keep $\mathcal Y$. \emph{Integrality.} $e_c$ is a sum over $g\in\prod\GL_{r_O}(\F_w)$ of the coefficients of the idempotent,
which have the denominator $\#\prod\GL_{r_O}(\F_w)$, dividing $n_{\rm ch}$, times a root of unity and the translate of the form by an integral point; so $n_{\rm ch}e_c$ keeps the forms integral above $\ell$. \emph{The representations.}
For an eigensystem of $\mathcal Y$ of auxiliary level, (b) gives the type. At a place $w\in Q$ the form has the parahoric type of $O$: an integral matrix whose reduction is block triangular acts through $\psi_w^{-1}$ of the determinant of
the block on the embedded indices, $\psi_w$ being the tame part of $\vartheta$ there, as $\omega$ is trivial at $Q$. By \cref{thm:B-galois}(b), (c), inertia at $w$ acts through $1$ and the tame character of $\psi_w^{-1}$, so $R$ is of
finite type there and $(X-\vartheta_\ell^{-1}(s))(X-1)$ kills each inertia element $s$ (\cref{lem:avatars}(c) for the avatar of $\vartheta^{-1}$ on inertia); conjugate duality carries this to the conjugate places. The normalising
character is trivial on inertia away from $\ell$, so the relation holds for the normalised representation.
\end{proof}

\subsection{Main summands of the lifts}\label{ss:main-summands-lifts}
A representation $T$ of rank $r_O$ is a \emph{main summand} for a descent pair $O$ and a system of Hecke eigenvalues $z$ if it is like a main summand and, for every tame twist $\vartheta$ of $O$ and every eigenform that lifts
$z$, the normalised representation $u$ attached to the lift satisfies
\[
  \operatorname{tr}u=\vartheta^{-1}_\ell\cdot\operatorname{tr}T+\sum_{e<t_O}\mathrm{cyc}^{-e}.
\]
\lean{New.Selmer.cmSign_star, New.Selmer.abs_descentPair_weight_add_liftWeight_le, New.Selmer.le_natAbs_descentPair_jump_and_le, New.Selmer.countP_low_weights_descentPair, New.Selmer.isMainSummandWeights_descentPair, New.Selmer.abs_mainPair_weight_add_liftWeight_le, New.Selmer.isMainSummandWeights_mainPair, New.Selmer.apply_one_eq_of_heckePolyEval_eq, New.Selmer.IsRepOfType.isMainSummandLike_twistByChar, New.Selmer.neg_coeff_eq_of_hasLiftHeckePolynomial, New.Selmer.eventually_trace_frob_eq_of_hasLiftHeckePolynomial, New.Selmer.exists_isMainSummand_of_isRepOfType, New.Selmer.map_jumps_sub_eq_map_targetParameter, New.Data.isEigenform_of_hasHeckePolyAt, New.Selmer.exists_isMainSummand_descentPair}
For the jumps of the descent pairs write $u_q(\sigma)=\lfloor\mathrm t/2\rfloor-\epsilon_\sigma(\lambda_q+w_j)$, $q<\mathrm r$, with the target parameter $\lambda$ at the place of $\sigma$ (\cref{prop:descent-pair-fits}); the gap $\Delta$ and the lift weights $w_j$
are those of \cref{prop:grids}.

\begin{lemma}[The jumps of the lifts]\label{lem:lift-jumps}
\begin{enumerate}[label=\textup{(\alph*)}]
\item For $q+1<k$ the jump $u_q$ is high, $\abs{\epsilon_\sigma u_q+w_j}\le\mathrm t$; for $k\le q+1\le\mathrm r$ it is low, $\Delta/4\le\abs{u_q}\le2m\Delta$. If $2m\Delta\le B$ and $B+\mathrm t<w_j$, exactly $m$ of the $u_q(\sigma)$ have modulus at most $B$.
\item For a main pair $(\nu,z)$ for $j$, the jump $z-1-\epsilon_\sigma c$ is high.
\item Let $y\ne0$ be a form of a descent pair $V$ of ranks $r<n$ that lifts the eigensystem $b$, and $w$ split off the places of $V$ and their conjugates. Then minus the coefficient of $X^{n-1}$ in the Hecke polynomial of $y$ at
$w$ is $q^t\nu\nu'^{-1}b(w,1)+\bigl(\sum_{i<t}q^i\bigr)\nu$, with $t=n-r$ and $\nu$, $\nu'$ the Frobenius values of the characters of $V$ on $W$ and on $V$.
\item In the setting of (c), for a pair $O$ whose character $\mu$ has $V$ as a twist by a tame $\vartheta$, with $r-t=2a$: at most Frobenius elements $g$, minus that coefficient times the normalising avatar of $V$ at $g$ is
$\theta_0(g)\theta(g)b(w,1)+\sum_{e<t}\mathrm{cyc}(g)^{-e}$, where $\theta_0$, $\theta$ are the avatars of $\vartheta^{-1}(x/\bar x)$ and of $\mu^{-1}(x/\bar x)\abs x^a$.
\end{enumerate}
\end{lemma}

\begin{proof}
(a) \emph{High jumps.} $\epsilon_\sigma u_q+w_j=\epsilon_\sigma\lfloor\mathrm t/2\rfloor-\lambda_q$ and $\abs{\lambda_q}\le k/2$ for $q+1<k$, with $k\le\mathrm t$; so the modulus is at most $\mathrm t$. \emph{Low jumps.} With $i=q+2-k\in[1,m]$ the parameter is
$\lambda_q+w_j=(j-i+\frac12)\Delta$, of modulus in $[\Delta/2,m\Delta]$, and $\lfloor\mathrm t/2\rfloor\le\Delta/4$ (\cref{prop:grids}(b)); hence $\Delta/4\le\abs{u_q}\le m\Delta+\Delta/4\le2m\Delta$. \emph{Counting.} The $m$ low jumps have modulus at most $2m\Delta\le B$; a high
jump has $\abs{u_q}\ge w_j-\mathrm t>B$.

(b) $\epsilon_\sigma(z-1-\epsilon_\sigma c)+w_j=\epsilon_\sigma(z-1)+(w_j-c)$, and $0\le z-1\le\mathrm r-1$, $0\le w_j-c\le\mathrm t-\mathrm r$; the modulus is at most $\mathrm t-1$.

(c) $y$ is spherical at $w$ with Hecke eigenvalues $s_i$, and, being nonzero, has the Hecke polynomial of a lift of $b$: $\prod_{l<t}(1-q^l\nu X)$ times the Hecke polynomial of rank $r$ of $b$ at $q^t\nu\nu'^{-1}X$
(\cref{prop:lift-descent}(c)). Comparing the linear terms of $\sum_i(-X)^iq^{i(i-1)/2}s_i$ gives $s_1=q^t\nu\nu'^{-1}b(w,1)+\bigl(\sum_{i<t}q^i\bigr)\nu$, the linear term of $\prod_{l<t}(1-q^l\nu X)$ being $-\sum_{i<t}q^i\nu$ and the
constant term of the polynomial of $b$ being $1$; and minus the coefficient of $X^{n-1}$ in the Hecke polynomial $\sum_j(-1)^jq^{j(j-1)/2}s_jX^{n-j}$ of $y$ (\cref{ss:B-galois}) is $s_1$.

(d) At a Frobenius element at $w\nmid\ell$, off the finitely many places where a character is ramified, the avatars take the values of the characters at a uniformiser (\cref{thm:B-galois}(a)), and the values of $\theta_0$ and $\theta$
multiply to $q^t/\nu'$; the normalising avatar takes the value $\nu^{-1}$ (its $d$-th power for the $d$-th power of Frobenius). Hence the normalised expression of (c) is $\theta_0\theta(g)b(w,1)+\sum_{i<t}q^i$, and $q^i=\mathrm{cyc}(g)^{-i}$
by \cref{lem:avatars}(b).
\end{proof}


\begin{theorem}[Main summands exist]\label{thm:main-summands-exist}
\begin{enumerate}[label=\textup{(\alph*)}]
\item The jumps $\lfloor\mathrm t/2\rfloor-\epsilon_\sigma(\lambda_q+w_j)$, $q<\mathrm r$, of the $j$-th descent pair are main summand weights of rank $\mathrm r$; the single jump $z-1-\epsilon_\sigma c$ of a main pair $(\nu,z)$ for $j$ is a main
summand weight of rank $1$ and string length $2z-1$.
\item A representation of type $c$ of rank $r_O$, empty level and main summand weights, twisted by the avatar of $\mu^{-1}(x/\bar x)\abs x^{a}$ with $r_O-t_O=2a$, is like a main summand.
\item For every Hecke eigensystem $z$ of a good form on $W_m$ and every $1\le j\le m$ there is a main summand for the $j$-th descent pair of the trivial twist and $z$; for every main pair there is a main summand of
rank one for its lift from $\U(1)$ and the constant system.
\item (Splitting.) The space of lifts for a twist $\vartheta$ and $j$ is spanned by forms in the image of the idempotent of a tuple of labels in lift setting, whose Hecke polynomials at the split places off the level
are those of a lift of one of finitely many eigensystems $z$ of good forms on $W_m$, independent of $\vartheta$.
\end{enumerate}
\end{theorem}

\begin{proof}
(a) \emph{Descent pairs.} By \cref{lem:lift-jumps}(a) each $u_q$ is low or high. Take $B=2(\Delta/2)^2$; then $2m\Delta\le B$ and $B+\mathrm t<w_j$ by the sizes of the data (\cref{prop:grids}), so exactly $m$ jumps at $\sigma$ are small, and
$m\le\mathrm rm/k$ as $k\le\mathrm r$. At $\sigma$ and $\bar\sigma$ the signs $\epsilon$ are opposite, so $u_q(\sigma)+u_q(\bar\sigma)=2\lfloor\mathrm t/2\rfloor\ge0$ and at most one of the two is negative: at most $\mathrm r$ negative jumps at the pair. The string length is
$\mathrm t$, as allowed. \emph{Main pairs.} By \cref{lem:lift-jumps}(b) the single jump is high, hence not small (it has modulus at least $w_j-\mathrm t>2(\Delta/2)^2$), and the jumps at $\sigma$ and $\bar\sigma$ add up to $2(z-1)\ge0$, so at most one is
negative; the string length is $2z-1\le\mathrm t$.

(b) \emph{Twisting.} The tame part of $\mu^{-1}$ is a power of $\omega$, trivial off the chosen places, so by \cref{lem:lift-settings}(e) the twist is a representation of a type whose jumps are lowered by the infinity type of the character;
these are the given main summand weights. \emph{Ramification above $p_0p_1$.} At a chosen place $w\nmid\ell$ an inertia element $s$ acts through the roots of the scaled labels raised to its tame exponent; these are roots of unity
of order dividing $p_i^{r_O!}-1$ and some $p_i^d-1$ with $d\le\mathrm r$, and $1$ is not among them for the tame generator, the scaled labels not having the root $1$. Hence the $N$-th powers of inertia act as the identity for
$N=p_i^{r_O!}-1$, a unit at $w$, and $\det(R\gamma-1)$ for the tame generator $\gamma$ is a product of $r_O$ differences $x-1$, each of valuation at most $v_\ell(p_i^d-1)\le v_\ell(n_{\rm ch})$ (\cref{lem:lift-settings}(b)). The places above
$p_0p_1$ that are not chosen are conjugates of chosen ones, and conjugate duality carries these statements there, with $(\gamma^c)^{-1}$ in place of $\gamma$. \emph{Large image.} Take a chosen place $w\nmid\ell$ with residue field of
order $q=p_i$ whose label is primitive with $r_O=\mathrm r$, or $X-1$ with $r_O=1$ (\cref{lem:avatars}(a), the lift setting). The roots $\xi$ of the scaled label have a common $e$-th power $y$ such that $y^x=1$ forces $q_2\mid jx$;
moreover $\xi^{q^{r_O}}=\xi$, and the powers $(\xi^{q^a})^{\ell^u}$, $a<r_O$, are distinct for every $u$. For $X-1$ the only root is the name of $\omega^{-j}$, of order $q_2$, and the last condition is empty. For a primitive label the
scaled label is again primitive, its roots are the Frobenius conjugates of a root $\eta$ of order $q^{\mathrm r}-1$, and $y$, their product, has order $q-1$, a multiple of $q_2^2$ by (D4) of \cref{def:data}; as $9\mid q-1$ and $\mathrm r$ is a
power of $3$, the conjugates of $\eta^{\ell^u}$ are distinct. Let $s$ be inertia of
tame exponent $1$ and $\tau$ the faithful character followed by $\iota$; then $\tau\xi$ is an eigenvalue of $R(s)$, and conjugation by a Frobenius $\phi$ raises $R(s)$ to the $q$-th power. The eigenvalues
$\tau\xi^{q^a}$ are pairwise incongruent modulo the maximal ideal, so $R(s)$ is regular semisimple residually, $R(\phi)$ permutes its eigenlines cyclically, and over all large coefficient rings the image of a lattice
spans all matrices, while $R(\phi^{r_O})$ is a scalar at an element of cyclotomic character $p_i^{r_O}$. \emph{A scalar of large order.} $X^e-\tau y$ kills $R(s)$, so $R(s^e)$ is the scalar $\tau y$. For $1\le t'\le4r_O+2$ and all $a$,
$q_2$ does not divide $j\ell^bt'$ for any $b$ (a prime $p\ne\ell$ dividing $q_0$ would give $p^s\mid jt'$, but $0<jt'<p^s$). Write the order of $y$ as $\ell^uk'$ with $\ell\nmid k'$. If $(y^{\ell^at'})^{\ell^u}=1$, then $q_2\mid j\ell^{a+u}t'$,
which is excluded; so $(y^{\ell^at'})^{\ell^u}$ and $1$ are distinct $k'$-th roots of unity, their difference has norm at least $\abs{k'}=1$ (\cref{lem:avatars}(b)), and it is an integral multiple of $y^{\ell^at'}-1$, which is therefore a
unit.

(c) \emph{Descent pairs.} Let $z\ne0$ be a good form with the system $p$. The idempotent $e_0$ fixes $z$ and is the sum of the idempotents of the good tuples, so some good tuple $c$ does not kill $z$. Its algebraic form is an
eigenform with the labels $c$: at a chosen place $\GL_{\mathrm r}$ of the local integers acts through the reduction, a parahoric type with one block. By \cref{lem:lift-settings}(a), (d), an attached $R$ is of a type for the $j$-th descent
pair of the trivial twist, and its traces are the limits of the eigenvalues $p(w,1)$ along Frobenius elements. Lower the jumps by $(\mathrm r-\mathrm t)/2-\epsilon_\sigma w_j$; by \cref{prop:descent-pair-fits} they become the $u_q(\sigma)$, main summand weights by (a).
Let $T$ be $R$ twisted by the avatar $\theta$ of $\mu^{-1}(x/\bar x)\abs x^a$; by (b) it is like a main summand. For every tame $\vartheta$ and every eigenform $y$ of a twisted pair lifting $p$, \cref{lem:lift-jumps}(d) gives the trace identity at
most Frobenius elements, $\theta_0$ being the avatar of $\vartheta^{-1}$; the trace of the normalised representation attached to $y$ is minus the coefficient of $X^{n-1}$ of its characteristic polynomial, the normalising avatar is
continuous, and the identity passes to all $g$ in the limit (\cref{thm:B-classfield}(d)). \emph{Main pairs.} The same argument for the lift from $\U(1)$, with the label $X-1$ (\cref{lem:lift-settings}(a)), the constant form, whose
system is that of the constant (\cref{prop:norm-L-value}(b)), and the jumps of (a) for main pairs.

(d) is \cref{prop:range-good-lift}, the Hecke polynomials of the lifts being those of \cref{thm:norm-one-lift}(d).
\end{proof}
